% SPDX-License-Identifier: GPL-3.0-or-later
% Copyright (C) 2026 Alessandro Greco
\chapter{Потоки, SENDME и управление перегрузкой}\label{cap:11-stream}
\aggiornatoal{23 сентября 2026 г.}{коммит torspec \texttt{8e1f245}; tor 0.4.9.13; предложения 289, 324}

\begin{obiettivi}
К концу этой главы вы сможете:
\begin{itemize}[nosep]
  \item открывать, использовать и закрывать TCP-поток через цепь с помощью BEGIN, CONNECTED, DATA и END;
  \item объяснять, почему разрешение имён выполняется на выходном реле и как работает RESOLVE;
  \item описывать исходное управление потоком Tor, включая окна цепи и потока и сообщения SENDME;
  \item понимать, почему SENDME аутентифицируются и что они доказывают;
  \item описывать управление перегрузкой из предложения 324 и роли XON и XOFF;
  \item загружать файл через собственноручно построенную трёхпереходную цепь и измерять окна реального выходного реле (лабораторная работа~5).
\end{itemize}
\end{obiettivi}

Цепь из \cref{cap:10-cifratura-relay} --- зашифрованный канал до выходного реле, но пока по нему не передаётся ничего полезного. Теперь задействуем его для TCP-соединений с Интернетом, которые Tor называет \textbf{потоками}. Затем рассмотрим проблему, которую должна решать каждая сеть и которая в Tor принимает особую форму: как не дать отправителю, передающему быстрее, чем сеть или получатель способны принимать, заполнить очереди реле.

% ---------------------------------------------------------------------------
\section{Потоки}\label{sec:stream}

Поток --- это TCP-соединение клиента, например с веб-сервером, переносимое через цепь. Одна цепь может одновременно нести множество потоков. Они различаются по \textbf{StreamID} в сообщениях RELAY (\cref{sec:relay}), который выбирает клиент и который должен быть ненулевым. Жизненный цикл потока показан на \cref{fig:stream}.

\begin{figure}[htbp]
\centering
% SPDX-License-Identifier: GPL-3.0-or-later
% Copyright (C) 2026 Alessandro Greco
%% Русское альтернативное описание EPUB.
\altfig{Схема потока клиент — выходной узел — сервер: BEGIN для example.com:80, возможные опережающие DATA до CONNECTED, отдельное разрешение DNS выходным узлом, TCP SYN/SYN-ACK с сервером, затем CONNECTED. DATA передаются по зашифрованной цепи Tor и по внешнему соединению TCP. SENDME уровня цепи в классическом управлении потоком использует шаг 100 ячеек DATA, не 31. FIN может привести к END (DONE). HTTPS может дополнительно защищать приложение за пределами Tor. Последовательность схематична.}
% Жизненный цикл stream.
\begin{tikzpicture}[font=\sffamily\small,>=Stealth,
  msg/.style={font=\ttfamily\scriptsize,fill=white,inner sep=1pt},
  tcp/.style={-{Stealth[length=2.5mm]},thick,torgrigio}]
  \node[client] (c) at (0,0) {Клиент};
  \node[relay,fill=torverde!30] (e) at (6.5,0) {Выход};
  \node[server] (s) at (10.5,0) {Веб-сервер};
  \foreach \x in {0,6.5,10.5} \draw[torgrigio] (\x,-0.5) -- (\x,-7.3);
  \draw[cellaf] (0,-1.0) -- node[msg,above]{BEGIN "example.com:80"} (6.5,-1.0);
  \draw[cellaf,dashed] (0,-1.6) -- node[msg,above]{DATA (опережающие)} (6.5,-1.6);
  \draw[tcp] (6.5,-1.3) -- node[msg,above]{TCP SYN} (10.5,-1.3);
  \draw[tcp] (10.5,-2.1) -- node[msg,above]{SYN-ACK} (6.5,-2.1);
  \draw[risposta] (6.5,-2.6) -- node[msg,above]{CONNECTED (IP, TTL)} (0,-2.6);
  \draw[tcp] (6.5,-2.9) -- node[msg,above]{запрос} (10.5,-2.9);
  \draw[tcp] (10.5,-3.6) -- node[msg,above]{ответ} (6.5,-3.6);
  \draw[risposta] (6.5,-4.1) -- node[msg,above]{DATA, DATA, \ldots} (0,-4.1);
  \draw[cellaf] (0,-4.8) -- node[msg,above]{SENDME (цепь, 100 DATA)} (6.5,-4.8);
  \draw[risposta] (6.5,-5.5) -- node[msg,above]{DATA, \ldots} (0,-5.5);
  \draw[tcp] (10.5,-6.1) -- node[msg,above]{FIN} (6.5,-6.1);
  \draw[risposta] (6.5,-6.7) -- node[msg,above]{END (DONE)} (0,-6.7);
  \node[font=\scriptsize,text=torgrigio] at (3.25,-7.6) {внутри цепи (3 перехода, зашифровано)};
  \node[font=\scriptsize,text=torgrigio] at (8.5,-7.6) {обычное TCP-соединение};
\end{tikzpicture}

\caption{Жизненный цикл потока к веб-серверу. Сообщения между клиентом и выходным реле идут через цепь под послойным шифрованием; между выходным реле и сервером находится обычное TCP-соединение, незашифрованное, если приложение не использует TLS. Клиент может отправить данные до получения CONNECTED (\emph{опережающие данные (optimistic data)}).}
\label{fig:stream}
\end{figure}

\subsection{Открытие потока: BEGIN и CONNECTED}

Чтобы открыть поток, клиент отправляет выходному реле сообщение \texttt{BEGIN} (команда RELAY 1) с новым StreamID. Тело задаёт пункт назначения~\spec{tor-spec}{Opening streams and transmitting data}:
\begin{center}
\begin{tabular}{lll}
\toprule
Поле & Длина & Содержимое \\
\midrule
ADDRPORT & переменная & \texttt{address:port}, завершённая нулевым байтом \\
FLAGS & 4 (необязательно) & предпочтения IPv4 и IPv6 \\
\bottomrule
\end{tabular}
\end{center}
Адресом может быть DNS-имя, IPv4-адрес или IPv6-адрес в квадратных скобках. Два небольших правила показывают приоритеты дизайна: строка должна отправляться \textbf{в нижнем регистре}, чтобы различия в капитализации не помогали отличать клиентов, а FLAGS \textbf{должно быть опущено}, если равно нулю. Биты FLAGS указывают, принимает ли клиент IPv6 (бит 1), отказывается ли от IPv4 (бит 2) и предпочитает ли IPv6 (бит 3).

Получив \texttt{BEGIN}, выходное реле:
\begin{enumerate}
  \item проверяет, что не является первым переходом цепи: канал, по которому пришла цепь, должен быть открыт аутентифицированным реле из консенсуса. Tor не должен работать как одношаговый прокси, что сделало бы выходные реле более привлекательными целями; такая попытка отклоняется \cella{DESTROY};
  \item применяет свою \emph{политику выхода (exit policy)}, список адресов и портов, к которым разрешены соединения (\cref{cap:22-relay});
  \item при необходимости разрешает имя и открывает TCP-соединение.
\end{enumerate}
Если всё успешно, реле отвечает \texttt{CONNECTED} (команда RELAY 4). Тело обычно содержит достигнутый IPv4-адрес (4 байта) и число секунд, в течение которых клиент может его кэшировать (TTL, 4 байта). Формат IPv6 немного отличается; разрешено и пустое тело. Клиент \emph{может} игнорировать адрес; если он его кэширует, то не должен повторно использовать адрес в других цепях. Нечестное выходное реле может солгать или вернуть уникальный адрес, чтобы впоследствии узнавать клиента.

\subsection{Передача данных: DATA}

Прикладные данные передаются сообщениями \texttt{DATA} (команда RELAY 2), каждое до 498 байт в традиционном формате. Выходное реле копирует их содержимое в TCP-соединение и обратно. Клиент может отправлять \texttt{DATA} сразу после \texttt{BEGIN}, не ожидая \texttt{CONNECTED}: выходное реле ставит эти сообщения в очередь и пересылает после открытия соединения. Такие \emph{опережающие данные (optimistic data)} экономят полный цикл прохождения туда и обратно по цепи, который в Tor легко может занимать несколько сотен миллисекунд.

\subsection{Закрытие потока: END}

Сторона, закрывающая соединение или обнаружившая ошибку, отправляет \texttt{END} (команда RELAY 3) с однобайтовой причиной~\spec{tor-spec}{Closing streams}. Наиболее частые причины: \texttt{DONE} (6, обычное закрытие), \texttt{CONNECTREFUSED} (3), \texttt{RESOLVEFAILED} (2), \texttt{TIMEOUT} (7) и \texttt{EXITPOLICY} (4); после последней следуют отклонённый адрес и TTL. Для созданных им потоков клиент должен, как правило, отправлять только \texttt{MISC} (1): как и с \cella{DESTROY}, нет причин сообщать другой стороне больше необходимого.

\subsection{DNS на выходном реле}\label{sec:stream-dns}

В \texttt{BEGIN} клиент может указать \emph{имя}, а не адрес: разрешение происходит на выходном реле. Это не просто удобство, а решение безопасности. Если бы клиент разрешал имена сам, он обращался бы к собственному DNS-резолверу вне Tor для каждого сайта, который собирается посетить. Наблюдатель таких запросов, начиная с провайдера подключения, узнал бы пункты назначения. Это знакомая \enquote{утечка DNS (DNS leak)}, демонстрируемая в упражнении~\ref{es:19-socks} и наблюдаемая с точки зрения tor в лабораторной работе~18 (\cref{sec:lab-opsec}).

Если клиенту нужен только адрес без открытия соединения, он использует \texttt{RESOLVE} (команда RELAY 11) с нуль-терминированным именем. Выходное реле отвечает \texttt{RESOLVED} (12), содержащим один или несколько ответов в формате тип, длина, значение и TTL: тип \texttt{0x04} для IPv4, \texttt{0x06} для IPv6, \texttt{0x00} для имени при обратном поиске и \texttt{0xF0}/\texttt{0xF1} для временной и постоянной ошибок~\spec{tor-spec}{Remote hostname lookup}. На этом механизме основан DNS-порт, который клиент tor может предоставлять локальным приложениям (опция \opt{DNSPort}).

\subsection{Потоки каталогов: BEGIN\_DIR}

\file{BEGIN\_DIR} (команда RELAY 13) открывает поток не в Интернет, а к собственному сервису каталогов реле. Так клиенты загружают консенсусы и дескрипторы (\cref{cap:12-directory}), не открывая незашифрованных соединений. Тело пустое, а политика выхода не применяется, поскольку поток не выходит из процесса tor.

% ---------------------------------------------------------------------------
\section{Исходное управление потоком}\label{sec:flusso-classico}

\subsection{Проблема}

Представим выходное реле с быстрым интернет-соединением и медленного клиента. Если выходное реле читает с сервера и без ограничений отправляет всё в цепь, данные накапливаются в очередях промежуточных реле, расходуя память и замедляя другие цепи. TCP решает эту проблему на каждой отдельной линии, но цепь --- последовательность разных TCP-соединений: обратное давление TCP (backpressure) между промежуточным реле и клиентом не доходит до выходного реле. Нужен механизм \emph{от конца до конца (end-to-end)}.

\subsection{Окна и SENDME}

Исходный механизм Tor использует \textbf{окна}, измеряемые в ячейках \texttt{DATA}~\spec{tor-spec}{Flow control}:
\begin{itemize}
  \item На уровне \textbf{цепи} отправитель имеет \emph{окно отправки (package window)}, начинающееся с 1000 ячеек или значения параметра консенсуса \texttt{circwindow}. Каждая отправленная \texttt{DATA} уменьшает его; при нуле отправка останавливается. Получатель ведёт симметричное \emph{окно доставки (deliver window)} и на каждые 100 полученных ячеек отправляет \texttt{SENDME} (команда RELAY 5) со StreamID 0, восстанавливая 100 ячеек в окне отправителя.
  \item На уровне \textbf{потока} окна начинаются с 500 ячеек и увеличиваются на 50 для каждого \texttt{SENDME} со StreamID этого потока.
\end{itemize}
Поскольку клиент может обмениваться данными с любым переходом (\cref{sec:tor1}), он поддерживает пару окон для каждого перехода. Эти окна считают только ячейки, созданные и потреблённые узлом, но не пересылаемые.

Следовательно, в любой момент в полёте может находиться не более 1000 ячеек на цепь и 500 на поток: примерно 500\,KB и 250\,KB. Это ограничивает использование памяти реле, но одновременно пропускную способность. Цепь с RTT в полсекунды не может превысить 1000 ячеек за полсекунды, то есть примерно 1\,MB/s, независимо от доступной полосы. Если сеть перегружена, даже 1000 ячеек в полёте слишком много: очереди и задержка растут.

\subsection{Аутентифицированные SENDME}

Вредоносный клиент мог бы отправлять \texttt{SENDME}, ничего не получив. Выходное реле продолжало бы отправлять данные, которые накапливались бы в сети, позволяя клиенту исчерпывать память реле или усиливать трафик с малой собственной стоимостью. Начиная с версии 0.4.1.1-alpha сообщения \texttt{SENDME} уровня цепи \textbf{аутентифицируются} (\proposta{289}). Их тело содержит версию (1), длину (20) и \textbf{накопительный дайджест} ячейки \texttt{DATA}, вызвавшей \texttt{SENDME}. Получатель SENDME знает этот digest, потому что вычислял его при отправке ячейки, и может сравнить значения. Несовпадение закрывает цепь.

Дайджест зависит от содержимого всех полученных ячеек, но клиент мог бы попытаться предсказать его, если заранее знает их содержимое, например при загрузке известного файла. Поэтому хотя бы одна ячейка в каждой группе из 100 должна содержать случайные байты: это одна из причин случайного заполнения сообщений RELAY (\cref{sec:relay}).

\begin{approfondimento}[Какой именно дайджест?]
Спецификация говорит о \enquote{дайджест ячейки перед \texttt{SENDME}}. В коде tor это полный 20-байтовый дайджест SHA-1 обратного потока после последней распознанной ячейки, сохранённый в поле \file{sendme\_digest} функцией \file{tor1\_relay\_digest\_matches\_v0} (\file{src/core/crypto/relay\_crypto\_tor1.c}). В CGO вместо него используется тег последней ячейки (\cref{sec:cgo}).
\end{approfondimento}

% ---------------------------------------------------------------------------
\section{Управление перегрузкой}\label{sec:cc}

Фиксированные окна имеют две противоположные слабости: они слишком малы для быстрых цепей и слишком велики для перегруженных. \proposta{324} Майка Перри ввела настоящее \textbf{управление перегрузкой} Tor, аналогичное TCP и принятое начиная с tor 0.4.7.7 (апрель 2022 года).

\subsection{Сигнал: время туда и обратно (RTT)}

Как выходное реле распознаёт перегрузку цепи? Используя \texttt{SENDME} как часы. Выходное реле записывает время отправки ячейки \texttt{DATA}, которая должна вызвать \texttt{SENDME}; после получения SENDME вычисляется \textbf{время туда и обратно} (RTT) цепи. Если RTT растёт относительно наблюдавшегося минимума, очереди вдоль цепи заполняются и отправителю следует замедлиться. Отправитель регулирует \emph{окно перегрузки} (\emph{cwnd}), стараясь удерживать его близко к произведению пропускной способности на задержку цепи: количеству данных, которое цепь способна переносить без создания очередей.

\subsection{Алгоритмы}

Предложение описывает три алгоритма, вдохновлённые TCP: Tor Westwood, Tor Vegas и Tor NOLA. После оценки в симуляторе Shadow (\cref{cap:19-rete-privata}) и в живой сети tor по умолчанию использует \textbf{Vegas} (константа \file{CC\_ALG\_DFLT} в \file{src/core/or/congestion\_control\_common.c}). Как TCP Vegas, он оценивает число \enquote{стоящих в очереди} ячеек, сравнивая текущее окно с окном, необходимым при минимальном наблюдавшемся RTT, затем увеличивает или уменьшает окно, удерживая оценку в заданных пределах.

\subsection{Параметры: почему 31?}

Управление перегрузкой согласуется отдельно для каждого перехода посредством ntor-v3 (\cref{sec:ntor3}). Клиент запрашивает его расширением \file{CC\_FIELD\_REQUEST}, а реле отвечает \file{CC\_FIELD\_RESPONSE} с \file{sendme\_inc}: числом ячеек \texttt{DATA}, после которого требуется \texttt{SENDME}. В \cref{cap:09-circuiti} выходное реле лабораторной сети вернул 31. Число не случайно: код tor определяет его как число ячеек, помещающихся в TLS-запись максимального размера,
\[ \file{TLS\_RECORD\_MAX\_CELLS} = \left\lfloor \frac{16384}{514} \right\rfloor = 31, \]
а начальное окно равно $4 \times 31 = 124$ ячейкам. Более частые SENDME, раз в 31 ячейку вместо 100, дают более частые измерения RTT и более быструю реакцию.

\subsection{Больше никаких окон потока: XON и XOFF}

При управлении перегрузкой сообщения \texttt{SENDME} уровня потока исчезают. Окна потока создавали дилемму: маленькие ограничивают скорость каждой загрузки; большие позволяют вредоносному клиенту заставить выходное реле буферизовать целое окно данных для каждого из множества открытых потоков. Предложение заменяет их двумя сообщениями:
\begin{description}
  \item[XOFF] (команда RELAY 43): \enquote{прекратить отправку по этому потоку}, посылается, когда очередь к приложению превышает порог;
  \item[XON] (команда RELAY 44): \enquote{возобновить}, с 4-байтовым полем скорости опустошения очереди, чтобы другая сторона могла адаптироваться. Поле называется \texttt{kbps\_ewma}, а \proposta{324} не определяет его единицу. Код tor вычисляет и интерпретирует его в единицах \textbf{1000 байт в секунду}, а не 1000 бит (\file{src/core/or/congestion\_control\_flow.c}): имя вводит в заблуждение, и буквальная интерпретация дала бы восьмикратную ошибку.
\end{description}
Сообщения \texttt{SENDME} уровня цепи, напротив, отправляются сразу после расшифрования и распознавания ячеек, до доставки приложению. Они измеряют сеть, а не приложение.

\begin{attenzione}[Управление перегрузкой и анонимность]
Любой механизм, реагирующий на трафик, может стать побочным каналом. Наблюдение за поведением окна способно раскрывать информацию о цепи. \proposta{324} посвящает таким рискам и мерам защиты отдельный раздел, включая ограничение частоты XON и XOFF. Это та же предосторожность, из-за которой в \cref{cap:08-celle} рекомендуется всегда отправлять причину NONE в \cella{DESTROY}: любая наблюдаемая вариация, зависящая от внутреннего состояния, выдаёт информацию. Атаки по признакам трафика рассматриваются в \cref{cap:26-attacchi}.
\end{attenzione}

% ---------------------------------------------------------------------------
\section{Лабораторная работа: загрузка файла через Tor}\label{sec:lab-stream}

\begin{laboratorio}[5]{Потоки и управление потоком}
\textbf{Цель.} Написать \file{labs/ru/код/torlab/circuit.py}, объединяющий предыдущие лабораторные работы в класс \texttt{Circuit}, и \file{labs/ru/код/torlab/stream.py}, открывающий потоки и передающий данные. Загрузить файл с локального веб-сервера по трёхпереходной цепи и измерить окна выходном реле.\\
\textbf{Предварительные условия.} Лабораторные работы~1--4; работающая лабораторная сеть; локальный веб-сервер (\cmd{python3 -m http.server 8080}).\\
\textbf{Оценочное время.} 3 часа.
\end{laboratorio}

\paragraph{Шаг 1: цепь одной строкой.} Функция \texttt{build} открывает канал к входному реле, создаёт первый переход и расширяет цепь, запрашивая управление перегрузкой только у последнего перехода, как делает tor:
\begin{lstlisting}[language=Python]
circ = build([relay("003r"), relay("004r"), relay("006e")])
print(circ.exit_extensions)             # [(2, b'\x1f')]: sendme_inc = 31
\end{lstlisting}

\paragraph{Шаг 2: открытие потока.}
\begin{lstlisting}[language=Python]
def begin_body(address, port, flag=0):
    body = f"{address.lower()}:{port}".encode() + b"\x00"
    return body + (struct.pack(">I", flag) if flag else b"")

s = Stream(circ, 1)
s.open("127.0.0.1", 8080)       # BEGIN, затем ожидание CONNECTED или END
\end{lstlisting}
Выходное реле отвечает \texttt{CONNECTED} с 8-байтовым телом \texttt{7f000001 00000708}: адрес 127.0.0.1 и TTL \texttt{0x708} = 1800 секунд.

\paragraph{Шаг 3: SENDME.} Цепь считает каждую ячейку \texttt{DATA}, полученную от выходного реле, и отправляет аутентифицированный \texttt{SENDME} через каждые \file{sendme\_inc} ячеек:
\begin{lstlisting}[language=Python]
def count_received_data(self):
    est = dict(self.exit_extensions)
    sendme_inc = est[2][0] if 2 in est else 100
    self.data_to_ack += 1
    if self.data_to_ack == sendme_inc:
        self.data_to_ack = 0
        digest = self.crypto.hops[-1].backward.digest.digest()
        body = struct.pack(">BH", 1, 20) + digest
        self.send(RelayMessage(RelayCommand.SENDME, 0, body))
\end{lstlisting}

\begin{attenzione}[Поучительная ошибка]
В первой версии лабораторной работы счётчик был атрибутом \emph{потока}. Первый поток, страница из двух ячеек, работал; второй, файл 1\,MB, прерывался выходным реле ячейкой \cella{DESTROY}. Причина: сообщения \texttt{SENDME} уровня цепи считают \texttt{DATA} из \emph{всех} потоков цепи. Выходное реле уже учло две ячейки первого потока и ожидало \texttt{SENDME} после 29-й ячейки второго. Мы отправляли его после 31-й с дайджестом, отличавшимся от ожидаемого, и выходное реле сделало именно то, что требует \proposta{289}: закрыло цепь. Перенос счётчика в класс \texttt{Circuit} исправил проблему.
\end{attenzione}

\paragraph{Шаг 4: загрузка.} Отправляем HTTP-запрос и читаем до \texttt{END}:
\begin{lstlisting}[language=Python]
s.send(b"GET /grande.bin HTTP/1.0\r\nHost: 127.0.0.1\r\n\r\n")
response = s.receive_all()
\end{lstlisting}
Для случайного файла размером \num{1048576} байт в одном из наших запусков:
\begin{lstlisting}[numbers=none]
stream 1: 209 байт за 0.05 с, DATA-ячеек 2, circuit SENDME пока 0,
          прочее ['END DONE']
stream 2: 1048781 байт за 0.18 с, DATA-ячеек 2120, circuit SENDME пока 68, прочее ['END DONE']
\end{lstlisting}
SHA-256-отпечаток полученного файла совпадает с исходным. 68 сообщений \texttt{SENDME} --- ровно $\lfloor (2 + 2120) / 31 \rfloor$: один счётчик на цепь. Ячеек 2120, а не минимальные $\lceil 1048781 / 498 \rceil = 2106$, поскольку выходное реле не ждёт накопления 498 байт перед отправкой ячейки, а упаковывает данные по мере чтения с сервера.

\paragraph{Шаг 5: наблюдение за окнами.} Теперь проведём эксперимент, невозможный для обычного клиента: загрузим тот же файл, \emph{никогда не отправляя} \texttt{SENDME}, и посчитаем, сколько ячеек выходное реле отправит до остановки. Повторим с управлением перегрузкой и без него:
\begin{lstlisting}[language=Python]
circ = build(path, cc=cc)
circ.count_received_data = lambda: None      # SENDME не отправляются
\end{lstlisting}
\begin{lstlisting}[numbers=none]
управление перегрузкой=True:  выходное реле останавливается после 124 ячеек DATA
управление перегрузкой=False: выходное реле останавливается после 500 ячеек DATA
\end{lstlisting}
С управлением перегрузкой выходное реле останавливается на 124 ячейках: это начальное окно перегрузки, $4 \times 31$. Без него выходное реле останавливается на 500, а не 1000: \emph{окно потока} (500) исчерпывается раньше окна цепи (1000), поскольку при исходном управлении потоком клиент также должен отправлять сообщения \texttt{SENDME} уровня потока. В обоих случаях выходное реле просто ждёт: данные остаются в его собственном буфере, а не в сети.

% ---------------------------------------------------------------------------
\section*{Итоги}
\begin{itemize}
  \item Поток открывается \texttt{BEGIN} (\texttt{address:port} в нижнем регистре), подтверждается \texttt{CONNECTED}, переносит \texttt{DATA} и закрывается \texttt{END}; клиент может отправлять данные до \texttt{CONNECTED}.
  \item Разрешение имён происходит на выходном реле, чтобы не раскрывать пункты назначения локальному DNS-резолверу; \texttt{RESOLVE} и \texttt{RESOLVED} позволяют разрешать имена без соединения.
  \item Исходное управление потоком использует окна 1000 ячеек на цепь и 500 на поток, с \texttt{SENDME} каждые 100 и 50 ячеек соответственно.
  \item Сообщения \texttt{SENDME} уровня цепи аутентифицируются накопительным дайджестом вызвавшей их ячейки (\proposta{289}).
  \item При управлении перегрузкой (\proposta{324}, алгоритм по умолчанию Vegas) окно адаптируется к RTT, измеренному через \texttt{SENDME}; \file{sendme\_inc} равен 31, числу ячеек в одном TLS-записи, а начальное окно равно 124.
  \item XON и XOFF заменяют окна потока.
\end{itemize}

% ---------------------------------------------------------------------------
\section*{Упражнения}

\begin{esercizio}[1]\label{es:11-begin}
Запишите в шестнадцатеричном виде тело сообщения \texttt{BEGIN} к \texttt{Example.COM} на порт 443 для клиента, который также принимает IPv6.
\end{esercizio}

\begin{esercizio}[1]\label{es:11-connected}
Интерпретируйте тело \texttt{CONNECTED} \texttt{c0000201 0000012c}. Что клиент должен сделать с этой информацией?
\end{esercizio}

\begin{esercizio}[1]\label{es:11-optimistic}
Оцените время, сэкономленное благодаря \emph{опережающей передаче данных (optimistic data)} для цепи с RTT 600\,ms и страницы, требующей одного HTTP-запроса. Почему выходное реле должно отбросить поставленные в очередь данные, если соединение не открылось?
\end{esercizio}

\begin{esercizio}[2]\label{es:11-velocita}
Какова максимальная скорость одного потока при исходном управлении потоком в цепи с RTT 400\,ms? А цепи с несколькими потоками? Сравните эти величины с начальным окном управления перегрузкой.
\end{esercizio}

\begin{esercizio}[2]\label{es:11-sendme-falso}
Опишите атаку на исчерпание памяти выходного реле, основанную на неаутентифицированных \texttt{SENDME}. Почему случайное заполнение в ячейках \texttt{DATA} необходимо для работы аутентификации?
\end{esercizio}

\begin{esercizio}[2]\label{es:11-resolve}
Расширьте \file{stream.py} функцией \texttt{resolve(name)}, использующей \texttt{RESOLVE} и интерпретирующей \texttt{RESOLVED}. Проверьте её в лабораторной сети с \texttt{localhost}. Как выходное реле отвечает на запрос несуществующего имени?
\end{esercizio}

\begin{esercizio}[3]\label{es:11-policy}
Попробуйте открыть поток на порт, который выходное реле лабораторной сети не принимает, после изменения политики выхода (\file{torrc} для \texttt{test006e}, затем \cmd{chutney hup}). Какую причину сообщает \texttt{END}? Интерпретируйте дополнительные поля.
\end{esercizio}

\begin{esercizio}[3]\label{es:11-rtt}
Измените лабораторную работу так, чтобы на клиенте измерять время между отправкой каждого \texttt{SENDME} и получением следующей ячейки \texttt{DATA}. Что именно вы измеряете? Сравните значения при наличии и отсутствии второй параллельной загрузки по той же цепи.
\end{esercizio}

% ---------------------------------------------------------------------------
\section*{Дополнительная литература}
\begin{itemize}
  \item \titolo{Tor Protocol Specification}, разделы \enquote{Opening streams and transmitting data}, \enquote{Closing streams}, \enquote{Remote hostname lookup} и \enquote{Flow control}~\cite{torspec}.
  \item Предложение 289, \titolo{Authenticating sendme cells to mitigate bandwidth attacks}.
  \item Предложение 324, \titolo{RTT-based Congestion Control for Tor}: мотивация, алгоритмы, параметры и анализ безопасности.
\end{itemize}
